\chapter*{Resumen}
\thispagestyle{empty}

% TODO (tesis): reescribir en voz ejecutada y sustituir "como hipótesis se
% plantea... <10%" por el error relativo medio REAL obtenido. Igual en el
% Abstract. Actualizar las categorías al esquema por firma de energía.

El crecimiento del IoT ha impulsado una demanda creciente de
dispositivos que operan con presupuestos de potencia muy limitados, donde la
eficiencia del software es tan determinante como la del hardware. Sin embargo,
los procesadores RISC-V de bajo consumo orientados a estas aplicaciones no
ofrecen ningún mecanismo que permita al firmware relacionar su ejecución con una
estimación de energía, obligando a los desarrolladores a recurrir a
equipos externos para caracterizar el comportamiento de sus
aplicaciones.

Este trabajo propone integrar directamente en el procesador CV32E40P, sobre la
plataforma PULPissimo, un módulo hardware que clasifica las instrucciones
ejecutadas y mantiene contadores accesibles desde el propio firmware, sin
requerir sistema operativo. La estimación se apoya en un modelo por categoría
de instrucción cuyos coeficientes se caracterizan experimentalmente mediante un
circuito de medición de potencia dedicado para la plataforma FPGA. La
instrumentación externa se utiliza para caracterizar y validar el modelo; una
vez obtenidos los coeficientes, el sistema puede estimar el consumo asociado a
la ejecución del firmware a partir de los contadores internos, sin repetir una
medición eléctrica completa para cada análisis. Como hipótesis de trabajo, se
plantea alcanzar un error relativo medio inferior al 10\,\% respecto a la
medición eléctrica de referencia.

\bigskip
\textbf{Palabras clave:} \scriptKeywords

\clearpage
\chapter*{Abstract}
\thispagestyle{empty}

The growth of the Internet of Things has driven increasing demand for devices
that operate under strict power budgets, where software efficiency is as
decisive as hardware efficiency. However, low-power RISC-V processors targeting
these applications provide no mechanism for firmware to relate its execution to
an energy estimate, forcing developers to rely on external laboratory equipment
to characterize the behavior of their applications.

This work proposes integrating directly into the CV32E40P processor, on the
PULPissimo platform, a hardware module that classifies executed instructions and
maintains counters accessible from firmware itself, without requiring an
operating system. The estimation relies on an instruction-category model whose
coefficients are experimentally characterized using external power measurement
instrumentation for the FPGA platform. This instrumentation is required for
characterization and validation; once the coefficients are obtained, the system
can estimate the energy consumption associated with firmware execution from the
internal counters without repeating a full electrical measurement for every
analysis. The working hypothesis targets a mean relative error below~10\,\%
with respect to the reference electrical measurement.

\bigskip
\textbf{Keywords:} RISC-V, CV32E40P, energy consumption, instruction classifier,
FPGA, PULPissimo, bare-metal, power measurement, CSR

\cleardoublepage

%%% Local Variables:
%%% mode: latex
%%% TeX-master: "main"
%%% End:
